\documentclass[12pt]{article}
\usepackage{amsmath,amssymb,amsfonts,amsthm}
\usepackage[margin=1in]{geometry}

\begin{document}

\begin{center}
{\Large\bfseries Chapter 4, Problem 6}\\[6pt]
Byron \& Fuller, \textit{Mathematics of Classical and Quantum Physics}
\end{center}

\bigskip

\textbf{Problem.} Either prove or find a counterexample to the following statements:
\begin{itemize}
\item[(a)] If $A$ and $B$ are $n \times n$ matrices, then $AB = 0$ implies that either $A = 0$ or $B = 0$.
\item[(b)] An $n \times n$ Hermitian matrix which is completely degenerate (all eigenvalues equal) is necessarily diagonal.
\item[(c)] If an $n \times n$ matrix is Hermitian, then all its powers are Hermitian.
\item[(d)] $\det(A + B) = \det A + \det B$.
\item[(e)] If $A$ and $B$ are $n \times n$ Hermitian matrices, then $AB$ is Hermitian.
\item[(f)] $\operatorname{tr} AB = \operatorname{tr} A \, \operatorname{tr} B$ in a finite-dimensional space.
\end{itemize}

\bigskip
\textbf{Solution.}

\textbf{(a) False.} Counterexample: Let
\[
A = \begin{pmatrix} 1 & 0 \\ 0 & 0 \end{pmatrix}, \quad B = \begin{pmatrix} 0 & 0 \\ 1 & 0 \end{pmatrix}.
\]
Then
\[
AB = \begin{pmatrix} 1 & 0 \\ 0 & 0 \end{pmatrix}\begin{pmatrix} 0 & 0 \\ 1 & 0 \end{pmatrix} = \begin{pmatrix} 0 & 0 \\ 0 & 0 \end{pmatrix} = 0,
\]
but neither $A = 0$ nor $B = 0$.

\bigskip
\textbf{(b) False.} Counterexample: If all eigenvalues equal $\lambda$, then $A = \lambda I$ (since $A$ is Hermitian and thus diagonalizable, and all eigenvalues are $\lambda$). But $\lambda I$ \emph{is} diagonal.

Wait---let us reconsider. The matrix $\lambda I$ is indeed diagonal. In fact the statement is \textbf{True}: If $A$ is Hermitian with all eigenvalues equal to $\lambda$, then $A$ can be diagonalized as $A = U \lambda I U^\dagger = \lambda I$, which is diagonal. However, note that the statement says ``necessarily diagonal,'' which is true since $\lambda I$ is diagonal regardless of basis.

Actually, let's be more careful. The claim is that a completely degenerate Hermitian matrix is ``necessarily diagonal.'' Since $A$ is Hermitian, it is unitarily diagonalizable: $A = U D U^\dagger$ where $D = \lambda I$. Then $A = U(\lambda I)U^\dagger = \lambda U U^\dagger = \lambda I$. So $A = \lambda I$, which is indeed diagonal.

\textbf{The statement is True.}

\bigskip
\textbf{(c) True.} If $A$ is Hermitian ($A^\dagger = A$), then for any positive integer $n$:
\[
(A^n)^\dagger = (A^\dagger)^n = A^n.
\]
So $A^n$ is Hermitian.

\bigskip
\textbf{(d) False.} Counterexample: Let $A = B = I$ (the $2 \times 2$ identity). Then
\[
\det(A + B) = \det(2I) = 4, \quad \det A + \det B = 1 + 1 = 2.
\]
Since $4 \neq 2$, the statement is false.

\bigskip
\textbf{(e) False.} Counterexample: Let
\[
A = \begin{pmatrix} 1 & 0 \\ 0 & 0 \end{pmatrix}, \quad B = \begin{pmatrix} 0 & 1 \\ 1 & 0 \end{pmatrix}.
\]
Both are Hermitian. Then
\[
AB = \begin{pmatrix} 0 & 1 \\ 0 & 0 \end{pmatrix},
\]
and
\[
(AB)^\dagger = \begin{pmatrix} 0 & 0 \\ 1 & 0 \end{pmatrix} \neq AB.
\]
So $AB$ is not Hermitian in general. (Note: $AB$ is Hermitian if and only if $A$ and $B$ commute, since $(AB)^\dagger = B^\dagger A^\dagger = BA$.)

\bigskip
\textbf{(f) False.} Counterexample: Let
\[
A = \begin{pmatrix} 1 & 0 \\ 0 & 0 \end{pmatrix}, \quad B = \begin{pmatrix} 0 & 0 \\ 0 & 1 \end{pmatrix}.
\]
Then $AB = 0$, so $\operatorname{tr}(AB) = 0$, but $\operatorname{tr} A \cdot \operatorname{tr} B = 1 \cdot 1 = 1 \neq 0$.

\end{document}
